% ============================================================
% Per-Sub-Agent Close-Up Diagrams (a)–(e)
% Fragment: % ============================================================
% Per-Sub-Agent Close-Up Diagrams (a)–(e)
% Fragment: % ============================================================
% Per-Sub-Agent Close-Up Diagrams (a)–(e)
% Fragment: % ============================================================
% Per-Sub-Agent Close-Up Diagrams (a)–(e)
% Fragment: \input{figs/agent_closeups} in your main doc.
% Requires: agentbox, database, statearrow styles;
%   all icon commands; all agent colors; stateBlue, aliceblue.
% ============================================================

\begin{figure}[p]
\centering

% ============================================================
% Shared tikz styles for all sub-figures
% ============================================================
\tikzset{
    agentmain/.style={
        rectangle, rounded corners=8pt, draw=#1!60!black, thick,
        minimum width=6.0cm, minimum height=1.2cm, align=center,
        fill=#1, drop shadow={opacity=0.10}, font=\sffamily
    },
    toolbox/.style={
        rectangle, rounded corners=4pt, draw=gray!50, thick,
        minimum width=4.8cm, minimum height=0.9cm, align=left,
        fill=white, font=\sffamily\small, inner sep=5pt
    },
    resource/.style={
        cylinder, shape border rotate=90, draw=gray!40,
        minimum width=1.0cm, minimum height=1.2cm,
        font=\sffamily\tiny\bfseries, align=center, fill=gray!8
    },
    stateio/.style={
        rectangle, rounded corners=4pt, draw=blue!30, thick,
        minimum width=4.0cm, minimum height=0.9cm, align=left,
        fill=stateBlue, font=\sffamily\small, inner sep=5pt
    },
    iolabel/.style={font=\sffamily\scriptsize\bfseries, text=blue!40!gray},
    rwlabel/.style={font=\sffamily\scriptsize, inner sep=1pt},
    conn/.style={->, >=latex, thick, draw=gray!50},
    stateconn/.style={<->, >=latex, thick, draw=blue!30}
}

% =======================  (a) OmicMiningAgent  =======================
\begin{minipage}[t]{0.48\textwidth}
\centering
\resizebox{\textwidth}{!}{%
\begin{tikzpicture}[node distance=0.6cm, font=\sffamily]
    % Agent header
    \node[agentmain=agentOrange] (agent)
        {\iconDNA{orange!60!black} \enskip \textbf{OmicMiningAgent}};

    % Tool
    \node[toolbox, below=0.8cm of agent] (tool1)
        {\texttt{omic\_analysis()}\\
         {\scriptsize NER $\rightarrow$ dataset search $\rightarrow$ DE $\rightarrow$ enrichment}};

    % Resources
    \node[resource, left=0.8cm of tool1, yshift=-0.8cm] (db1) {TOSG\\DB};
    \node[resource, below=0.2cm of db1] (db2) {GLiNER\\NER};
    \node[resource, right=0.8cm of tool1, yshift=-0.8cm] (db3) {DESeq2\\/ R};
    \node[resource, below=0.2cm of db3] (db4) {KEGG\\Enrich};

    % State I/O
    \node[stateio, below=1.8cm of tool1] (reads)
        {\textcolor{gray!60!black}{\textbf{R:}} \texttt{query}};
    \node[stateio, below=0.3cm of reads]  (writes)
        {\textcolor{blue!50!black}{\textbf{W:}} \texttt{top\_genes[]}, \texttt{disease},\\
         \quad\; \texttt{cell\_type}, \texttt{pathways[]}};

    % Connections
    \draw[conn] (agent) -- (tool1);
    \draw[conn, draw=gray!30] (tool1.south) -- ++(0,-0.3) -| (db1.north);
    \draw[conn, draw=gray!30] (tool1.south) -- ++(0,-0.3) -| (db2.north);
    \draw[conn, draw=gray!30] (tool1.south) -- ++(0,-0.3) -| (db3.north);
    \draw[conn, draw=gray!30] (tool1.south) -- ++(0,-0.3) -| (db4.north);

    \node[iolabel, left=0.1cm of reads.north west, anchor=south east] {AgentState};
    \draw[stateconn] (tool1.south east) -- ++(0.6,0) |- (reads.east);
    \draw[stateconn] (tool1.south east) -- ++(0.6,0) |- (writes.east);
\end{tikzpicture}%
}
\par\smallskip{\small (a) OmicMiningAgent}
\end{minipage}%
\hfill
% =======================  (b) BioMarkerKGAgent  =======================
\begin{minipage}[t]{0.48\textwidth}
\centering
\resizebox{\textwidth}{!}{%
\begin{tikzpicture}[node distance=0.6cm, font=\sffamily]
    % Agent header
    \node[agentmain=agentCyan] (agent)
        {\iconGraph{cyan!60!black} \enskip \textbf{BioMarkerKGAgent}};

    % Tool
    \node[toolbox, below=0.8cm of agent] (tool1)
        {\texttt{query\_knowledge\_graph()}\\
         {\scriptsize GRetriever: GAT + LLM over PrimeKG}};

    % Resources
    \node[resource, left=0.8cm of tool1, yshift=-0.8cm] (db1) {Neo4j\\PrimeKG};
    \node[resource, below=0.2cm of db1] (db2) {GRetriever\\(PyG)};
    \node[resource, right=0.8cm of tool1, yshift=-0.8cm] (db3) {OpenAI\\Embed};
    \node[resource, below=0.2cm of db3] (db4) {PCST\\Subgraph};

    % State I/O
    \node[stateio, below=1.8cm of tool1] (reads)
        {\textcolor{gray!60!black}{\textbf{R:}} \texttt{top\_genes},
         \texttt{previous\_results}};
    \node[stateio, below=0.3cm of reads]  (writes)
        {\textcolor{blue!50!black}{\textbf{W:}} \texttt{agent\_outputs[]}\\
         {\scriptsize (pathway chains, GO terms, PPIs)}};

    % Connections
    \draw[conn] (agent) -- (tool1);
    \draw[conn, draw=gray!30] (tool1.south) -- ++(0,-0.3) -| (db1.north);
    \draw[conn, draw=gray!30] (tool1.south) -- ++(0,-0.3) -| (db2.north);
    \draw[conn, draw=gray!30] (tool1.south) -- ++(0,-0.3) -| (db3.north);
    \draw[conn, draw=gray!30] (tool1.south) -- ++(0,-0.3) -| (db4.north);

    \node[iolabel, left=0.1cm of reads.north west, anchor=south east] {AgentState};
    \draw[stateconn] (tool1.south east) -- ++(0.6,0) |- (reads.east);
    \draw[stateconn] (tool1.south east) -- ++(0.6,0) |- (writes.east);
\end{tikzpicture}%
}
\par\smallskip{\small (b) BioMarkerKGAgent}
\end{minipage}

\vspace{0.8cm}

% =======================  (c) PubMedResearcher  =======================
\begin{minipage}[t]{0.48\textwidth}
\centering
\resizebox{\textwidth}{!}{%
\begin{tikzpicture}[node distance=0.6cm, font=\sffamily]
    % Agent header
    \node[agentmain=agentPurple] (agent)
        {\iconSearch{violet} \enskip \textbf{PubMedResearcher}};

    % Tools (2)
    \node[toolbox, below=0.8cm of agent, xshift=-1.5cm] (tool1)
        {\texttt{search\_pubmed()}\\
         {\scriptsize Full-text PDF/XML download}};
    \node[toolbox, below=0.8cm of agent, xshift=1.5cm] (tool2)
        {\texttt{pubmed\_search\_fast()}\\
         {\scriptsize NCBI E-utilities (abstracts)}};

    % Resources
    \node[resource, below=1.0cm of tool1] (db1) {PubMed\\API};
    \node[resource, right=0.3cm of db1] (db2) {PMC\\PDFs};
    \node[resource, below=1.0cm of tool2] (db3) {DOI\\Cache};

    % State I/O
    \node[stateio, below=2.4cm of agent] (reads)
        {\textcolor{gray!60!black}{\textbf{R:}} \texttt{top\_genes},
         \texttt{previous\_results}};
    \node[stateio, below=0.3cm of reads]  (writes)
        {\textcolor{blue!50!black}{\textbf{W:}} \texttt{paper\_dois[]},
         \texttt{paper\_pmids[]}\\
         {\scriptsize (per-gene literature with citations)}};

    % Connections
    \draw[conn] (agent.south) -- ++(0,-0.3) -| (tool1.north);
    \draw[conn] (agent.south) -- ++(0,-0.3) -| (tool2.north);
    \draw[conn, draw=gray!30] (tool1) -- (db1);
    \draw[conn, draw=gray!30] (tool1) -- (db2);
    \draw[conn, draw=gray!30] (tool2) -- (db3);

    \node[iolabel, left=0.1cm of reads.north west, anchor=south east] {AgentState};
    \draw[stateconn] (reads.east) -- ++(0.6,0) |- ($(agent.south east)+(0.6,-0.3)$);
    \draw[stateconn] (writes.east) -- ++(0.6,0) |- ($(agent.south east)+(0.6,-0.3)$);
\end{tikzpicture}%
}
\par\smallskip{\small (c) PubMedResearcher}
\end{minipage}%
\hfill
% =======================  (d) GoogleSearcher  =======================
\begin{minipage}[t]{0.48\textwidth}
\centering
\resizebox{\textwidth}{!}{%
\begin{tikzpicture}[node distance=0.6cm, font=\sffamily]
    % Agent header
    \node[agentmain=googleGray] (agent)
        {\iconSearch{gray!60!black} \enskip \textbf{GoogleSearcher}};

    % Tool
    \node[toolbox, below=0.8cm of agent] (tool1)
        {\texttt{google\_search()}\\
         {\scriptsize Custom Search API + w3m extraction}};

    % Resources
    \node[resource, left=0.8cm of tool1, yshift=-0.8cm] (db1) {Google\\CSE};
    \node[resource, right=0.8cm of tool1, yshift=-0.8cm] (db2) {w3m\\Browser};
    \node[resource, below=0.2cm of db2] (db3) {LLM\\Filter};

    % State I/O
    \node[stateio, below=1.8cm of tool1] (reads)
        {\textcolor{gray!60!black}{\textbf{R:}} \texttt{previous\_results}
         (hypotheses)};
    \node[stateio, below=0.3cm of reads]  (writes)
        {\textcolor{blue!50!black}{\textbf{W:}} \texttt{agent\_outputs[]}\\
         {\scriptsize (clinical trials, translational viability)}};

    % Connections
    \draw[conn] (agent) -- (tool1);
    \draw[conn, draw=gray!30] (tool1.south) -- ++(0,-0.3) -| (db1.north);
    \draw[conn, draw=gray!30] (tool1.south) -- ++(0,-0.3) -| (db2.north);
    \draw[conn, draw=gray!30] (tool1.south) -- ++(0,-0.3) -| (db3.north);

    \node[iolabel, left=0.1cm of reads.north west, anchor=south east] {AgentState};
    \draw[stateconn] (tool1.south east) -- ++(0.6,0) |- (reads.east);
    \draw[stateconn] (tool1.south east) -- ++(0.6,0) |- (writes.east);
\end{tikzpicture}%
}
\par\smallskip{\small (d) GoogleSearcher}
\end{minipage}

\vspace{0.8cm}

% =======================  (e) ScientistsAgent  =======================
\begin{minipage}[t]{\textwidth}
\centering
\resizebox{0.85\textwidth}{!}{%
\begin{tikzpicture}[node distance=0.6cm, font=\sffamily]
    % Agent header (wider for compound agent)
    \node[agentmain=agentTeal, minimum width=8cm] (agent)
        {\iconBulb{teal!60!black} \enskip \textbf{ScientistsAgent}
         {\small\itshape (Multi-Expert PI)}};

    % Router
    \node[rectangle, rounded corners=4pt, draw=teal!50, thick,
          fill=teal!8, minimum width=5cm, minimum height=0.9cm,
          align=center, font=\sffamily\small,
          below=0.7cm of agent] (router)
        {\iconBrain{teal!60!black} \enskip \textbf{Expert Router (PI)}\\
         {\scriptsize Structured LLM routing $\rightarrow$ selects 2--3 experts}};

    % 4 Expert boxes in a row
    \node[toolbox, fill=yellow!8, below=0.9cm of router, xshift=-5.2cm,
          minimum width=3.6cm, minimum height=1.8cm] (exp1)
        {\textbf{GenomicsExpert}\\[2pt]
         \texttt{\scriptsize query\_genomics-}\\
         \texttt{\scriptsize expert\_kb()}\\
         {\scriptsize HyperRAG KB}};

    \node[toolbox, fill=pink!8, right=0.3cm of exp1,
          minimum width=3.6cm, minimum height=1.8cm] (exp2)
        {\textbf{NeuroscienceExpert}\\[2pt]
         \texttt{\scriptsize query\_neuroscience-}\\
         \texttt{\scriptsize expert\_kb()}\\
         {\scriptsize HyperRAG KB}};

    \node[toolbox, fill=lime!8, right=0.3cm of exp2,
          minimum width=3.6cm, minimum height=1.8cm] (exp3)
        {\textbf{LongevityBiostats}\\[2pt]
         \texttt{\scriptsize query\_longevity-}\\
         \texttt{\scriptsize biostatsexpert\_kb()}\\
         {\scriptsize HyperRAG KB}};

    \node[toolbox, fill=cyan!8, right=0.3cm of exp3,
          minimum width=3.6cm, minimum height=1.8cm] (exp4)
        {\textbf{BioinformaticsExpert}\\[2pt]
         \texttt{\scriptsize query\_bioinformatics-}\\
         \texttt{\scriptsize expert\_kb()}\\
         {\scriptsize HyperRAG KB}};

    % Synthesis
    \node[rectangle, rounded corners=4pt, draw=teal!50, thick,
          fill=teal!8, minimum width=5cm, minimum height=0.9cm,
          align=center, font=\sffamily\small,
          below=0.9cm of $(exp2.south)!0.5!(exp3.south)$] (synth)
        {\iconRobot{teal!60!black} \enskip \textbf{PI Synthesis}\\
         {\scriptsize Integrates expert perspectives $\rightarrow$ ranked hypotheses}};

    % State I/O (below synthesis)
    \node[stateio, minimum width=6cm, below=0.8cm of synth, xshift=-2.5cm] (reads)
        {\textcolor{gray!60!black}{\textbf{R:}} \texttt{ALL previous\_results},
         \texttt{top\_genes}, \texttt{paper\_dois}};
    \node[stateio, minimum width=6cm, below=0.3cm of reads]  (writes)
        {\textcolor{blue!50!black}{\textbf{W:}} \texttt{agent\_outputs[]}
         + \texttt{expert\_routing\{\}}\\
         {\scriptsize (ranked hypotheses, validation plans, expert contributions)}};

    % RAG resources
    \node[resource, below=0.15cm of exp1] (rag1) {Author\\KB};
    \node[resource, below=0.15cm of exp2] (rag2) {Author\\KB};
    \node[resource, below=0.15cm of exp3] (rag3) {Author\\KB};
    \node[resource, below=0.15cm of exp4] (rag4) {Author\\KB};

    % Connections: agent -> router
    \draw[conn] (agent) -- (router);

    % Router -> experts (fan-out)
    \draw[conn, draw=teal!40] (router.south) -- ++(0,-0.25) -| (exp1.north);
    \draw[conn, draw=teal!40] (router.south) -- ++(0,-0.25) -| (exp2.north);
    \draw[conn, draw=teal!40] (router.south) -- ++(0,-0.25) -| (exp3.north);
    \draw[conn, draw=teal!40] (router.south) -- ++(0,-0.25) -| (exp4.north);

    % Experts -> synthesis (fan-in)
    \draw[conn, draw=teal!40] (exp1.south) -- ++(0,-0.95) -| (synth.north west);
    \draw[conn, draw=teal!40] (exp2.south) -- ++(0,-0.95) -| ($(synth.north)!0.33!(synth.north west)$);
    \draw[conn, draw=teal!40] (exp3.south) -- ++(0,-0.95) -| ($(synth.north)!0.33!(synth.north east)$);
    \draw[conn, draw=teal!40] (exp4.south) -- ++(0,-0.95) -| (synth.north east);

    % State connections
    \node[iolabel, left=0.1cm of reads.north west, anchor=south east] {AgentState};
    \draw[stateconn] (synth.south east) -- ++(0.6,0) |- (reads.east);
    \draw[stateconn] (synth.south east) -- ++(0.6,0) |- (writes.east);
\end{tikzpicture}%
}
\par\smallskip{\small (e) ScientistsAgent}
\end{minipage}

\caption{Close-up views of each sub-agent in the OmniCellAgent system, showing
tools, external resources, and AgentState read/write interfaces.
\textbf{(a)~OmicMiningAgent} performs differential expression analysis via a
single \texttt{omic\_analysis} tool backed by the OmniCellTOSG single-cell
database, GLiNER NER for entity extraction, and R/DESeq2 for statistical
testing; it writes the foundational \texttt{top\_genes} list that seeds all
downstream agents.
\textbf{(b)~BioMarkerKGAgent} queries the PrimeKG knowledge graph through a
GRetriever model (GAT + LLM on Neo4j) to map genes to pathways, protein
interactions, and GO terms, contextualising the omics hits within known
biology.
\textbf{(c)~PubMedResearcher} has two tools---a full-text downloader
(\texttt{search\_pubmed}) and a fast abstract search
(\texttt{pubmed\_search\_fast})---producing per-gene literature reviews with
DOI/PMID citations that populate the shared state for traceability.
\textbf{(d)~GoogleSearcher} uses the Google Custom Search API with w3m-based
page extraction to assess clinical trial pipelines and translational viability
of the proposed targets.
\textbf{(e)~ScientistsAgent} is a compound multi-expert agent: a PI-level
router selects 2--3 domain experts (Genomics, Neuroscience,
Longevity/Biostatistics, Bioinformatics), each optionally backed by a
HyperRAG knowledge base built from curated author publications; a synthesis
node integrates their perspectives into ranked mechanistic hypotheses with
proposed validation experiments.}
\label{fig:agent-closeups}
\end{figure}
 in your main doc.
% Requires: agentbox, database, statearrow styles;
%   all icon commands; all agent colors; stateBlue, aliceblue.
% ============================================================

\begin{figure}[p]
\centering

% ============================================================
% Shared tikz styles for all sub-figures
% ============================================================
\tikzset{
    agentmain/.style={
        rectangle, rounded corners=8pt, draw=#1!60!black, thick,
        minimum width=6.0cm, minimum height=1.2cm, align=center,
        fill=#1, drop shadow={opacity=0.10}, font=\sffamily
    },
    toolbox/.style={
        rectangle, rounded corners=4pt, draw=gray!50, thick,
        minimum width=4.8cm, minimum height=0.9cm, align=left,
        fill=white, font=\sffamily\small, inner sep=5pt
    },
    resource/.style={
        cylinder, shape border rotate=90, draw=gray!40,
        minimum width=1.0cm, minimum height=1.2cm,
        font=\sffamily\tiny\bfseries, align=center, fill=gray!8
    },
    stateio/.style={
        rectangle, rounded corners=4pt, draw=blue!30, thick,
        minimum width=4.0cm, minimum height=0.9cm, align=left,
        fill=stateBlue, font=\sffamily\small, inner sep=5pt
    },
    iolabel/.style={font=\sffamily\scriptsize\bfseries, text=blue!40!gray},
    rwlabel/.style={font=\sffamily\scriptsize, inner sep=1pt},
    conn/.style={->, >=latex, thick, draw=gray!50},
    stateconn/.style={<->, >=latex, thick, draw=blue!30}
}

% =======================  (a) OmicMiningAgent  =======================
\begin{minipage}[t]{0.48\textwidth}
\centering
\resizebox{\textwidth}{!}{%
\begin{tikzpicture}[node distance=0.6cm, font=\sffamily]
    % Agent header
    \node[agentmain=agentOrange] (agent)
        {\iconDNA{orange!60!black} \enskip \textbf{OmicMiningAgent}};

    % Tool
    \node[toolbox, below=0.8cm of agent] (tool1)
        {\texttt{omic\_analysis()}\\
         {\scriptsize NER $\rightarrow$ dataset search $\rightarrow$ DE $\rightarrow$ enrichment}};

    % Resources
    \node[resource, left=0.8cm of tool1, yshift=-0.8cm] (db1) {TOSG\\DB};
    \node[resource, below=0.2cm of db1] (db2) {GLiNER\\NER};
    \node[resource, right=0.8cm of tool1, yshift=-0.8cm] (db3) {DESeq2\\/ R};
    \node[resource, below=0.2cm of db3] (db4) {KEGG\\Enrich};

    % State I/O
    \node[stateio, below=1.8cm of tool1] (reads)
        {\textcolor{gray!60!black}{\textbf{R:}} \texttt{query}};
    \node[stateio, below=0.3cm of reads]  (writes)
        {\textcolor{blue!50!black}{\textbf{W:}} \texttt{top\_genes[]}, \texttt{disease},\\
         \quad\; \texttt{cell\_type}, \texttt{pathways[]}};

    % Connections
    \draw[conn] (agent) -- (tool1);
    \draw[conn, draw=gray!30] (tool1.south) -- ++(0,-0.3) -| (db1.north);
    \draw[conn, draw=gray!30] (tool1.south) -- ++(0,-0.3) -| (db2.north);
    \draw[conn, draw=gray!30] (tool1.south) -- ++(0,-0.3) -| (db3.north);
    \draw[conn, draw=gray!30] (tool1.south) -- ++(0,-0.3) -| (db4.north);

    \node[iolabel, left=0.1cm of reads.north west, anchor=south east] {AgentState};
    \draw[stateconn] (tool1.south east) -- ++(0.6,0) |- (reads.east);
    \draw[stateconn] (tool1.south east) -- ++(0.6,0) |- (writes.east);
\end{tikzpicture}%
}
\par\smallskip{\small (a) OmicMiningAgent}
\end{minipage}%
\hfill
% =======================  (b) BioMarkerKGAgent  =======================
\begin{minipage}[t]{0.48\textwidth}
\centering
\resizebox{\textwidth}{!}{%
\begin{tikzpicture}[node distance=0.6cm, font=\sffamily]
    % Agent header
    \node[agentmain=agentCyan] (agent)
        {\iconGraph{cyan!60!black} \enskip \textbf{BioMarkerKGAgent}};

    % Tool
    \node[toolbox, below=0.8cm of agent] (tool1)
        {\texttt{query\_knowledge\_graph()}\\
         {\scriptsize GRetriever: GAT + LLM over PrimeKG}};

    % Resources
    \node[resource, left=0.8cm of tool1, yshift=-0.8cm] (db1) {Neo4j\\PrimeKG};
    \node[resource, below=0.2cm of db1] (db2) {GRetriever\\(PyG)};
    \node[resource, right=0.8cm of tool1, yshift=-0.8cm] (db3) {OpenAI\\Embed};
    \node[resource, below=0.2cm of db3] (db4) {PCST\\Subgraph};

    % State I/O
    \node[stateio, below=1.8cm of tool1] (reads)
        {\textcolor{gray!60!black}{\textbf{R:}} \texttt{top\_genes},
         \texttt{previous\_results}};
    \node[stateio, below=0.3cm of reads]  (writes)
        {\textcolor{blue!50!black}{\textbf{W:}} \texttt{agent\_outputs[]}\\
         {\scriptsize (pathway chains, GO terms, PPIs)}};

    % Connections
    \draw[conn] (agent) -- (tool1);
    \draw[conn, draw=gray!30] (tool1.south) -- ++(0,-0.3) -| (db1.north);
    \draw[conn, draw=gray!30] (tool1.south) -- ++(0,-0.3) -| (db2.north);
    \draw[conn, draw=gray!30] (tool1.south) -- ++(0,-0.3) -| (db3.north);
    \draw[conn, draw=gray!30] (tool1.south) -- ++(0,-0.3) -| (db4.north);

    \node[iolabel, left=0.1cm of reads.north west, anchor=south east] {AgentState};
    \draw[stateconn] (tool1.south east) -- ++(0.6,0) |- (reads.east);
    \draw[stateconn] (tool1.south east) -- ++(0.6,0) |- (writes.east);
\end{tikzpicture}%
}
\par\smallskip{\small (b) BioMarkerKGAgent}
\end{minipage}

\vspace{0.8cm}

% =======================  (c) PubMedResearcher  =======================
\begin{minipage}[t]{0.48\textwidth}
\centering
\resizebox{\textwidth}{!}{%
\begin{tikzpicture}[node distance=0.6cm, font=\sffamily]
    % Agent header
    \node[agentmain=agentPurple] (agent)
        {\iconSearch{violet} \enskip \textbf{PubMedResearcher}};

    % Tools (2)
    \node[toolbox, below=0.8cm of agent, xshift=-1.5cm] (tool1)
        {\texttt{search\_pubmed()}\\
         {\scriptsize Full-text PDF/XML download}};
    \node[toolbox, below=0.8cm of agent, xshift=1.5cm] (tool2)
        {\texttt{pubmed\_search\_fast()}\\
         {\scriptsize NCBI E-utilities (abstracts)}};

    % Resources
    \node[resource, below=1.0cm of tool1] (db1) {PubMed\\API};
    \node[resource, right=0.3cm of db1] (db2) {PMC\\PDFs};
    \node[resource, below=1.0cm of tool2] (db3) {DOI\\Cache};

    % State I/O
    \node[stateio, below=2.4cm of agent] (reads)
        {\textcolor{gray!60!black}{\textbf{R:}} \texttt{top\_genes},
         \texttt{previous\_results}};
    \node[stateio, below=0.3cm of reads]  (writes)
        {\textcolor{blue!50!black}{\textbf{W:}} \texttt{paper\_dois[]},
         \texttt{paper\_pmids[]}\\
         {\scriptsize (per-gene literature with citations)}};

    % Connections
    \draw[conn] (agent.south) -- ++(0,-0.3) -| (tool1.north);
    \draw[conn] (agent.south) -- ++(0,-0.3) -| (tool2.north);
    \draw[conn, draw=gray!30] (tool1) -- (db1);
    \draw[conn, draw=gray!30] (tool1) -- (db2);
    \draw[conn, draw=gray!30] (tool2) -- (db3);

    \node[iolabel, left=0.1cm of reads.north west, anchor=south east] {AgentState};
    \draw[stateconn] (reads.east) -- ++(0.6,0) |- ($(agent.south east)+(0.6,-0.3)$);
    \draw[stateconn] (writes.east) -- ++(0.6,0) |- ($(agent.south east)+(0.6,-0.3)$);
\end{tikzpicture}%
}
\par\smallskip{\small (c) PubMedResearcher}
\end{minipage}%
\hfill
% =======================  (d) GoogleSearcher  =======================
\begin{minipage}[t]{0.48\textwidth}
\centering
\resizebox{\textwidth}{!}{%
\begin{tikzpicture}[node distance=0.6cm, font=\sffamily]
    % Agent header
    \node[agentmain=googleGray] (agent)
        {\iconSearch{gray!60!black} \enskip \textbf{GoogleSearcher}};

    % Tool
    \node[toolbox, below=0.8cm of agent] (tool1)
        {\texttt{google\_search()}\\
         {\scriptsize Custom Search API + w3m extraction}};

    % Resources
    \node[resource, left=0.8cm of tool1, yshift=-0.8cm] (db1) {Google\\CSE};
    \node[resource, right=0.8cm of tool1, yshift=-0.8cm] (db2) {w3m\\Browser};
    \node[resource, below=0.2cm of db2] (db3) {LLM\\Filter};

    % State I/O
    \node[stateio, below=1.8cm of tool1] (reads)
        {\textcolor{gray!60!black}{\textbf{R:}} \texttt{previous\_results}
         (hypotheses)};
    \node[stateio, below=0.3cm of reads]  (writes)
        {\textcolor{blue!50!black}{\textbf{W:}} \texttt{agent\_outputs[]}\\
         {\scriptsize (clinical trials, translational viability)}};

    % Connections
    \draw[conn] (agent) -- (tool1);
    \draw[conn, draw=gray!30] (tool1.south) -- ++(0,-0.3) -| (db1.north);
    \draw[conn, draw=gray!30] (tool1.south) -- ++(0,-0.3) -| (db2.north);
    \draw[conn, draw=gray!30] (tool1.south) -- ++(0,-0.3) -| (db3.north);

    \node[iolabel, left=0.1cm of reads.north west, anchor=south east] {AgentState};
    \draw[stateconn] (tool1.south east) -- ++(0.6,0) |- (reads.east);
    \draw[stateconn] (tool1.south east) -- ++(0.6,0) |- (writes.east);
\end{tikzpicture}%
}
\par\smallskip{\small (d) GoogleSearcher}
\end{minipage}

\vspace{0.8cm}

% =======================  (e) ScientistsAgent  =======================
\begin{minipage}[t]{\textwidth}
\centering
\resizebox{0.85\textwidth}{!}{%
\begin{tikzpicture}[node distance=0.6cm, font=\sffamily]
    % Agent header (wider for compound agent)
    \node[agentmain=agentTeal, minimum width=8cm] (agent)
        {\iconBulb{teal!60!black} \enskip \textbf{ScientistsAgent}
         {\small\itshape (Multi-Expert PI)}};

    % Router
    \node[rectangle, rounded corners=4pt, draw=teal!50, thick,
          fill=teal!8, minimum width=5cm, minimum height=0.9cm,
          align=center, font=\sffamily\small,
          below=0.7cm of agent] (router)
        {\iconBrain{teal!60!black} \enskip \textbf{Expert Router (PI)}\\
         {\scriptsize Structured LLM routing $\rightarrow$ selects 2--3 experts}};

    % 4 Expert boxes in a row
    \node[toolbox, fill=yellow!8, below=0.9cm of router, xshift=-5.2cm,
          minimum width=3.6cm, minimum height=1.8cm] (exp1)
        {\textbf{GenomicsExpert}\\[2pt]
         \texttt{\scriptsize query\_genomics-}\\
         \texttt{\scriptsize expert\_kb()}\\
         {\scriptsize HyperRAG KB}};

    \node[toolbox, fill=pink!8, right=0.3cm of exp1,
          minimum width=3.6cm, minimum height=1.8cm] (exp2)
        {\textbf{NeuroscienceExpert}\\[2pt]
         \texttt{\scriptsize query\_neuroscience-}\\
         \texttt{\scriptsize expert\_kb()}\\
         {\scriptsize HyperRAG KB}};

    \node[toolbox, fill=lime!8, right=0.3cm of exp2,
          minimum width=3.6cm, minimum height=1.8cm] (exp3)
        {\textbf{LongevityBiostats}\\[2pt]
         \texttt{\scriptsize query\_longevity-}\\
         \texttt{\scriptsize biostatsexpert\_kb()}\\
         {\scriptsize HyperRAG KB}};

    \node[toolbox, fill=cyan!8, right=0.3cm of exp3,
          minimum width=3.6cm, minimum height=1.8cm] (exp4)
        {\textbf{BioinformaticsExpert}\\[2pt]
         \texttt{\scriptsize query\_bioinformatics-}\\
         \texttt{\scriptsize expert\_kb()}\\
         {\scriptsize HyperRAG KB}};

    % Synthesis
    \node[rectangle, rounded corners=4pt, draw=teal!50, thick,
          fill=teal!8, minimum width=5cm, minimum height=0.9cm,
          align=center, font=\sffamily\small,
          below=0.9cm of $(exp2.south)!0.5!(exp3.south)$] (synth)
        {\iconRobot{teal!60!black} \enskip \textbf{PI Synthesis}\\
         {\scriptsize Integrates expert perspectives $\rightarrow$ ranked hypotheses}};

    % State I/O (below synthesis)
    \node[stateio, minimum width=6cm, below=0.8cm of synth, xshift=-2.5cm] (reads)
        {\textcolor{gray!60!black}{\textbf{R:}} \texttt{ALL previous\_results},
         \texttt{top\_genes}, \texttt{paper\_dois}};
    \node[stateio, minimum width=6cm, below=0.3cm of reads]  (writes)
        {\textcolor{blue!50!black}{\textbf{W:}} \texttt{agent\_outputs[]}
         + \texttt{expert\_routing\{\}}\\
         {\scriptsize (ranked hypotheses, validation plans, expert contributions)}};

    % RAG resources
    \node[resource, below=0.15cm of exp1] (rag1) {Author\\KB};
    \node[resource, below=0.15cm of exp2] (rag2) {Author\\KB};
    \node[resource, below=0.15cm of exp3] (rag3) {Author\\KB};
    \node[resource, below=0.15cm of exp4] (rag4) {Author\\KB};

    % Connections: agent -> router
    \draw[conn] (agent) -- (router);

    % Router -> experts (fan-out)
    \draw[conn, draw=teal!40] (router.south) -- ++(0,-0.25) -| (exp1.north);
    \draw[conn, draw=teal!40] (router.south) -- ++(0,-0.25) -| (exp2.north);
    \draw[conn, draw=teal!40] (router.south) -- ++(0,-0.25) -| (exp3.north);
    \draw[conn, draw=teal!40] (router.south) -- ++(0,-0.25) -| (exp4.north);

    % Experts -> synthesis (fan-in)
    \draw[conn, draw=teal!40] (exp1.south) -- ++(0,-0.95) -| (synth.north west);
    \draw[conn, draw=teal!40] (exp2.south) -- ++(0,-0.95) -| ($(synth.north)!0.33!(synth.north west)$);
    \draw[conn, draw=teal!40] (exp3.south) -- ++(0,-0.95) -| ($(synth.north)!0.33!(synth.north east)$);
    \draw[conn, draw=teal!40] (exp4.south) -- ++(0,-0.95) -| (synth.north east);

    % State connections
    \node[iolabel, left=0.1cm of reads.north west, anchor=south east] {AgentState};
    \draw[stateconn] (synth.south east) -- ++(0.6,0) |- (reads.east);
    \draw[stateconn] (synth.south east) -- ++(0.6,0) |- (writes.east);
\end{tikzpicture}%
}
\par\smallskip{\small (e) ScientistsAgent}
\end{minipage}

\caption{Close-up views of each sub-agent in the OmniCellAgent system, showing
tools, external resources, and AgentState read/write interfaces.
\textbf{(a)~OmicMiningAgent} performs differential expression analysis via a
single \texttt{omic\_analysis} tool backed by the OmniCellTOSG single-cell
database, GLiNER NER for entity extraction, and R/DESeq2 for statistical
testing; it writes the foundational \texttt{top\_genes} list that seeds all
downstream agents.
\textbf{(b)~BioMarkerKGAgent} queries the PrimeKG knowledge graph through a
GRetriever model (GAT + LLM on Neo4j) to map genes to pathways, protein
interactions, and GO terms, contextualising the omics hits within known
biology.
\textbf{(c)~PubMedResearcher} has two tools---a full-text downloader
(\texttt{search\_pubmed}) and a fast abstract search
(\texttt{pubmed\_search\_fast})---producing per-gene literature reviews with
DOI/PMID citations that populate the shared state for traceability.
\textbf{(d)~GoogleSearcher} uses the Google Custom Search API with w3m-based
page extraction to assess clinical trial pipelines and translational viability
of the proposed targets.
\textbf{(e)~ScientistsAgent} is a compound multi-expert agent: a PI-level
router selects 2--3 domain experts (Genomics, Neuroscience,
Longevity/Biostatistics, Bioinformatics), each optionally backed by a
HyperRAG knowledge base built from curated author publications; a synthesis
node integrates their perspectives into ranked mechanistic hypotheses with
proposed validation experiments.}
\label{fig:agent-closeups}
\end{figure}
 in your main doc.
% Requires: agentbox, database, statearrow styles;
%   all icon commands; all agent colors; stateBlue, aliceblue.
% ============================================================

\begin{figure}[p]
\centering

% ============================================================
% Shared tikz styles for all sub-figures
% ============================================================
\tikzset{
    agentmain/.style={
        rectangle, rounded corners=8pt, draw=#1!60!black, thick,
        minimum width=6.0cm, minimum height=1.2cm, align=center,
        fill=#1, drop shadow={opacity=0.10}, font=\sffamily
    },
    toolbox/.style={
        rectangle, rounded corners=4pt, draw=gray!50, thick,
        minimum width=4.8cm, minimum height=0.9cm, align=left,
        fill=white, font=\sffamily\small, inner sep=5pt
    },
    resource/.style={
        cylinder, shape border rotate=90, draw=gray!40,
        minimum width=1.0cm, minimum height=1.2cm,
        font=\sffamily\tiny\bfseries, align=center, fill=gray!8
    },
    stateio/.style={
        rectangle, rounded corners=4pt, draw=blue!30, thick,
        minimum width=4.0cm, minimum height=0.9cm, align=left,
        fill=stateBlue, font=\sffamily\small, inner sep=5pt
    },
    iolabel/.style={font=\sffamily\scriptsize\bfseries, text=blue!40!gray},
    rwlabel/.style={font=\sffamily\scriptsize, inner sep=1pt},
    conn/.style={->, >=latex, thick, draw=gray!50},
    stateconn/.style={<->, >=latex, thick, draw=blue!30}
}

% =======================  (a) OmicMiningAgent  =======================
\begin{minipage}[t]{0.48\textwidth}
\centering
\resizebox{\textwidth}{!}{%
\begin{tikzpicture}[node distance=0.6cm, font=\sffamily]
    % Agent header
    \node[agentmain=agentOrange] (agent)
        {\iconDNA{orange!60!black} \enskip \textbf{OmicMiningAgent}};

    % Tool
    \node[toolbox, below=0.8cm of agent] (tool1)
        {\texttt{omic\_analysis()}\\
         {\scriptsize NER $\rightarrow$ dataset search $\rightarrow$ DE $\rightarrow$ enrichment}};

    % Resources
    \node[resource, left=0.8cm of tool1, yshift=-0.8cm] (db1) {TOSG\\DB};
    \node[resource, below=0.2cm of db1] (db2) {GLiNER\\NER};
    \node[resource, right=0.8cm of tool1, yshift=-0.8cm] (db3) {DESeq2\\/ R};
    \node[resource, below=0.2cm of db3] (db4) {KEGG\\Enrich};

    % State I/O
    \node[stateio, below=1.8cm of tool1] (reads)
        {\textcolor{gray!60!black}{\textbf{R:}} \texttt{query}};
    \node[stateio, below=0.3cm of reads]  (writes)
        {\textcolor{blue!50!black}{\textbf{W:}} \texttt{top\_genes[]}, \texttt{disease},\\
         \quad\; \texttt{cell\_type}, \texttt{pathways[]}};

    % Connections
    \draw[conn] (agent) -- (tool1);
    \draw[conn, draw=gray!30] (tool1.south) -- ++(0,-0.3) -| (db1.north);
    \draw[conn, draw=gray!30] (tool1.south) -- ++(0,-0.3) -| (db2.north);
    \draw[conn, draw=gray!30] (tool1.south) -- ++(0,-0.3) -| (db3.north);
    \draw[conn, draw=gray!30] (tool1.south) -- ++(0,-0.3) -| (db4.north);

    \node[iolabel, left=0.1cm of reads.north west, anchor=south east] {AgentState};
    \draw[stateconn] (tool1.south east) -- ++(0.6,0) |- (reads.east);
    \draw[stateconn] (tool1.south east) -- ++(0.6,0) |- (writes.east);
\end{tikzpicture}%
}
\par\smallskip{\small (a) OmicMiningAgent}
\end{minipage}%
\hfill
% =======================  (b) BioMarkerKGAgent  =======================
\begin{minipage}[t]{0.48\textwidth}
\centering
\resizebox{\textwidth}{!}{%
\begin{tikzpicture}[node distance=0.6cm, font=\sffamily]
    % Agent header
    \node[agentmain=agentCyan] (agent)
        {\iconGraph{cyan!60!black} \enskip \textbf{BioMarkerKGAgent}};

    % Tool
    \node[toolbox, below=0.8cm of agent] (tool1)
        {\texttt{query\_knowledge\_graph()}\\
         {\scriptsize GRetriever: GAT + LLM over PrimeKG}};

    % Resources
    \node[resource, left=0.8cm of tool1, yshift=-0.8cm] (db1) {Neo4j\\PrimeKG};
    \node[resource, below=0.2cm of db1] (db2) {GRetriever\\(PyG)};
    \node[resource, right=0.8cm of tool1, yshift=-0.8cm] (db3) {OpenAI\\Embed};
    \node[resource, below=0.2cm of db3] (db4) {PCST\\Subgraph};

    % State I/O
    \node[stateio, below=1.8cm of tool1] (reads)
        {\textcolor{gray!60!black}{\textbf{R:}} \texttt{top\_genes},
         \texttt{previous\_results}};
    \node[stateio, below=0.3cm of reads]  (writes)
        {\textcolor{blue!50!black}{\textbf{W:}} \texttt{agent\_outputs[]}\\
         {\scriptsize (pathway chains, GO terms, PPIs)}};

    % Connections
    \draw[conn] (agent) -- (tool1);
    \draw[conn, draw=gray!30] (tool1.south) -- ++(0,-0.3) -| (db1.north);
    \draw[conn, draw=gray!30] (tool1.south) -- ++(0,-0.3) -| (db2.north);
    \draw[conn, draw=gray!30] (tool1.south) -- ++(0,-0.3) -| (db3.north);
    \draw[conn, draw=gray!30] (tool1.south) -- ++(0,-0.3) -| (db4.north);

    \node[iolabel, left=0.1cm of reads.north west, anchor=south east] {AgentState};
    \draw[stateconn] (tool1.south east) -- ++(0.6,0) |- (reads.east);
    \draw[stateconn] (tool1.south east) -- ++(0.6,0) |- (writes.east);
\end{tikzpicture}%
}
\par\smallskip{\small (b) BioMarkerKGAgent}
\end{minipage}

\vspace{0.8cm}

% =======================  (c) PubMedResearcher  =======================
\begin{minipage}[t]{0.48\textwidth}
\centering
\resizebox{\textwidth}{!}{%
\begin{tikzpicture}[node distance=0.6cm, font=\sffamily]
    % Agent header
    \node[agentmain=agentPurple] (agent)
        {\iconSearch{violet} \enskip \textbf{PubMedResearcher}};

    % Tools (2)
    \node[toolbox, below=0.8cm of agent, xshift=-1.5cm] (tool1)
        {\texttt{search\_pubmed()}\\
         {\scriptsize Full-text PDF/XML download}};
    \node[toolbox, below=0.8cm of agent, xshift=1.5cm] (tool2)
        {\texttt{pubmed\_search\_fast()}\\
         {\scriptsize NCBI E-utilities (abstracts)}};

    % Resources
    \node[resource, below=1.0cm of tool1] (db1) {PubMed\\API};
    \node[resource, right=0.3cm of db1] (db2) {PMC\\PDFs};
    \node[resource, below=1.0cm of tool2] (db3) {DOI\\Cache};

    % State I/O
    \node[stateio, below=2.4cm of agent] (reads)
        {\textcolor{gray!60!black}{\textbf{R:}} \texttt{top\_genes},
         \texttt{previous\_results}};
    \node[stateio, below=0.3cm of reads]  (writes)
        {\textcolor{blue!50!black}{\textbf{W:}} \texttt{paper\_dois[]},
         \texttt{paper\_pmids[]}\\
         {\scriptsize (per-gene literature with citations)}};

    % Connections
    \draw[conn] (agent.south) -- ++(0,-0.3) -| (tool1.north);
    \draw[conn] (agent.south) -- ++(0,-0.3) -| (tool2.north);
    \draw[conn, draw=gray!30] (tool1) -- (db1);
    \draw[conn, draw=gray!30] (tool1) -- (db2);
    \draw[conn, draw=gray!30] (tool2) -- (db3);

    \node[iolabel, left=0.1cm of reads.north west, anchor=south east] {AgentState};
    \draw[stateconn] (reads.east) -- ++(0.6,0) |- ($(agent.south east)+(0.6,-0.3)$);
    \draw[stateconn] (writes.east) -- ++(0.6,0) |- ($(agent.south east)+(0.6,-0.3)$);
\end{tikzpicture}%
}
\par\smallskip{\small (c) PubMedResearcher}
\end{minipage}%
\hfill
% =======================  (d) GoogleSearcher  =======================
\begin{minipage}[t]{0.48\textwidth}
\centering
\resizebox{\textwidth}{!}{%
\begin{tikzpicture}[node distance=0.6cm, font=\sffamily]
    % Agent header
    \node[agentmain=googleGray] (agent)
        {\iconSearch{gray!60!black} \enskip \textbf{GoogleSearcher}};

    % Tool
    \node[toolbox, below=0.8cm of agent] (tool1)
        {\texttt{google\_search()}\\
         {\scriptsize Custom Search API + w3m extraction}};

    % Resources
    \node[resource, left=0.8cm of tool1, yshift=-0.8cm] (db1) {Google\\CSE};
    \node[resource, right=0.8cm of tool1, yshift=-0.8cm] (db2) {w3m\\Browser};
    \node[resource, below=0.2cm of db2] (db3) {LLM\\Filter};

    % State I/O
    \node[stateio, below=1.8cm of tool1] (reads)
        {\textcolor{gray!60!black}{\textbf{R:}} \texttt{previous\_results}
         (hypotheses)};
    \node[stateio, below=0.3cm of reads]  (writes)
        {\textcolor{blue!50!black}{\textbf{W:}} \texttt{agent\_outputs[]}\\
         {\scriptsize (clinical trials, translational viability)}};

    % Connections
    \draw[conn] (agent) -- (tool1);
    \draw[conn, draw=gray!30] (tool1.south) -- ++(0,-0.3) -| (db1.north);
    \draw[conn, draw=gray!30] (tool1.south) -- ++(0,-0.3) -| (db2.north);
    \draw[conn, draw=gray!30] (tool1.south) -- ++(0,-0.3) -| (db3.north);

    \node[iolabel, left=0.1cm of reads.north west, anchor=south east] {AgentState};
    \draw[stateconn] (tool1.south east) -- ++(0.6,0) |- (reads.east);
    \draw[stateconn] (tool1.south east) -- ++(0.6,0) |- (writes.east);
\end{tikzpicture}%
}
\par\smallskip{\small (d) GoogleSearcher}
\end{minipage}

\vspace{0.8cm}

% =======================  (e) ScientistsAgent  =======================
\begin{minipage}[t]{\textwidth}
\centering
\resizebox{0.85\textwidth}{!}{%
\begin{tikzpicture}[node distance=0.6cm, font=\sffamily]
    % Agent header (wider for compound agent)
    \node[agentmain=agentTeal, minimum width=8cm] (agent)
        {\iconBulb{teal!60!black} \enskip \textbf{ScientistsAgent}
         {\small\itshape (Multi-Expert PI)}};

    % Router
    \node[rectangle, rounded corners=4pt, draw=teal!50, thick,
          fill=teal!8, minimum width=5cm, minimum height=0.9cm,
          align=center, font=\sffamily\small,
          below=0.7cm of agent] (router)
        {\iconBrain{teal!60!black} \enskip \textbf{Expert Router (PI)}\\
         {\scriptsize Structured LLM routing $\rightarrow$ selects 2--3 experts}};

    % 4 Expert boxes in a row
    \node[toolbox, fill=yellow!8, below=0.9cm of router, xshift=-5.2cm,
          minimum width=3.6cm, minimum height=1.8cm] (exp1)
        {\textbf{GenomicsExpert}\\[2pt]
         \texttt{\scriptsize query\_genomics-}\\
         \texttt{\scriptsize expert\_kb()}\\
         {\scriptsize HyperRAG KB}};

    \node[toolbox, fill=pink!8, right=0.3cm of exp1,
          minimum width=3.6cm, minimum height=1.8cm] (exp2)
        {\textbf{NeuroscienceExpert}\\[2pt]
         \texttt{\scriptsize query\_neuroscience-}\\
         \texttt{\scriptsize expert\_kb()}\\
         {\scriptsize HyperRAG KB}};

    \node[toolbox, fill=lime!8, right=0.3cm of exp2,
          minimum width=3.6cm, minimum height=1.8cm] (exp3)
        {\textbf{LongevityBiostats}\\[2pt]
         \texttt{\scriptsize query\_longevity-}\\
         \texttt{\scriptsize biostatsexpert\_kb()}\\
         {\scriptsize HyperRAG KB}};

    \node[toolbox, fill=cyan!8, right=0.3cm of exp3,
          minimum width=3.6cm, minimum height=1.8cm] (exp4)
        {\textbf{BioinformaticsExpert}\\[2pt]
         \texttt{\scriptsize query\_bioinformatics-}\\
         \texttt{\scriptsize expert\_kb()}\\
         {\scriptsize HyperRAG KB}};

    % Synthesis
    \node[rectangle, rounded corners=4pt, draw=teal!50, thick,
          fill=teal!8, minimum width=5cm, minimum height=0.9cm,
          align=center, font=\sffamily\small,
          below=0.9cm of $(exp2.south)!0.5!(exp3.south)$] (synth)
        {\iconRobot{teal!60!black} \enskip \textbf{PI Synthesis}\\
         {\scriptsize Integrates expert perspectives $\rightarrow$ ranked hypotheses}};

    % State I/O (below synthesis)
    \node[stateio, minimum width=6cm, below=0.8cm of synth, xshift=-2.5cm] (reads)
        {\textcolor{gray!60!black}{\textbf{R:}} \texttt{ALL previous\_results},
         \texttt{top\_genes}, \texttt{paper\_dois}};
    \node[stateio, minimum width=6cm, below=0.3cm of reads]  (writes)
        {\textcolor{blue!50!black}{\textbf{W:}} \texttt{agent\_outputs[]}
         + \texttt{expert\_routing\{\}}\\
         {\scriptsize (ranked hypotheses, validation plans, expert contributions)}};

    % RAG resources
    \node[resource, below=0.15cm of exp1] (rag1) {Author\\KB};
    \node[resource, below=0.15cm of exp2] (rag2) {Author\\KB};
    \node[resource, below=0.15cm of exp3] (rag3) {Author\\KB};
    \node[resource, below=0.15cm of exp4] (rag4) {Author\\KB};

    % Connections: agent -> router
    \draw[conn] (agent) -- (router);

    % Router -> experts (fan-out)
    \draw[conn, draw=teal!40] (router.south) -- ++(0,-0.25) -| (exp1.north);
    \draw[conn, draw=teal!40] (router.south) -- ++(0,-0.25) -| (exp2.north);
    \draw[conn, draw=teal!40] (router.south) -- ++(0,-0.25) -| (exp3.north);
    \draw[conn, draw=teal!40] (router.south) -- ++(0,-0.25) -| (exp4.north);

    % Experts -> synthesis (fan-in)
    \draw[conn, draw=teal!40] (exp1.south) -- ++(0,-0.95) -| (synth.north west);
    \draw[conn, draw=teal!40] (exp2.south) -- ++(0,-0.95) -| ($(synth.north)!0.33!(synth.north west)$);
    \draw[conn, draw=teal!40] (exp3.south) -- ++(0,-0.95) -| ($(synth.north)!0.33!(synth.north east)$);
    \draw[conn, draw=teal!40] (exp4.south) -- ++(0,-0.95) -| (synth.north east);

    % State connections
    \node[iolabel, left=0.1cm of reads.north west, anchor=south east] {AgentState};
    \draw[stateconn] (synth.south east) -- ++(0.6,0) |- (reads.east);
    \draw[stateconn] (synth.south east) -- ++(0.6,0) |- (writes.east);
\end{tikzpicture}%
}
\par\smallskip{\small (e) ScientistsAgent}
\end{minipage}

\caption{Close-up views of each sub-agent in the OmniCellAgent system, showing
tools, external resources, and AgentState read/write interfaces.
\textbf{(a)~OmicMiningAgent} performs differential expression analysis via a
single \texttt{omic\_analysis} tool backed by the OmniCellTOSG single-cell
database, GLiNER NER for entity extraction, and R/DESeq2 for statistical
testing; it writes the foundational \texttt{top\_genes} list that seeds all
downstream agents.
\textbf{(b)~BioMarkerKGAgent} queries the PrimeKG knowledge graph through a
GRetriever model (GAT + LLM on Neo4j) to map genes to pathways, protein
interactions, and GO terms, contextualising the omics hits within known
biology.
\textbf{(c)~PubMedResearcher} has two tools---a full-text downloader
(\texttt{search\_pubmed}) and a fast abstract search
(\texttt{pubmed\_search\_fast})---producing per-gene literature reviews with
DOI/PMID citations that populate the shared state for traceability.
\textbf{(d)~GoogleSearcher} uses the Google Custom Search API with w3m-based
page extraction to assess clinical trial pipelines and translational viability
of the proposed targets.
\textbf{(e)~ScientistsAgent} is a compound multi-expert agent: a PI-level
router selects 2--3 domain experts (Genomics, Neuroscience,
Longevity/Biostatistics, Bioinformatics), each optionally backed by a
HyperRAG knowledge base built from curated author publications; a synthesis
node integrates their perspectives into ranked mechanistic hypotheses with
proposed validation experiments.}
\label{fig:agent-closeups}
\end{figure}
 in your main doc.
% Requires: agentbox, database, statearrow styles;
%   all icon commands; all agent colors; stateBlue, aliceblue.
% ============================================================

\begin{figure}[p]
\centering

% ============================================================
% Shared tikz styles for all sub-figures
% ============================================================
\tikzset{
    agentmain/.style={
        rectangle, rounded corners=8pt, draw=#1!60!black, thick,
        minimum width=6.0cm, minimum height=1.2cm, align=center,
        fill=#1, drop shadow={opacity=0.10}, font=\sffamily
    },
    toolbox/.style={
        rectangle, rounded corners=4pt, draw=gray!50, thick,
        minimum width=4.8cm, minimum height=0.9cm, align=left,
        fill=white, font=\sffamily\small, inner sep=5pt
    },
    resource/.style={
        cylinder, shape border rotate=90, draw=gray!40,
        minimum width=1.0cm, minimum height=1.2cm,
        font=\sffamily\tiny\bfseries, align=center, fill=gray!8
    },
    stateio/.style={
        rectangle, rounded corners=4pt, draw=blue!30, thick,
        minimum width=4.0cm, minimum height=0.9cm, align=left,
        fill=stateBlue, font=\sffamily\small, inner sep=5pt
    },
    iolabel/.style={font=\sffamily\scriptsize\bfseries, text=blue!40!gray},
    rwlabel/.style={font=\sffamily\scriptsize, inner sep=1pt},
    conn/.style={->, >=latex, thick, draw=gray!50},
    stateconn/.style={<->, >=latex, thick, draw=blue!30}
}

% =======================  (a) OmicMiningAgent  =======================
\begin{minipage}[t]{0.48\textwidth}
\centering
\resizebox{\textwidth}{!}{%
\begin{tikzpicture}[node distance=0.6cm, font=\sffamily]
    % Agent header
    \node[agentmain=agentOrange] (agent)
        {\iconDNA{orange!60!black} \enskip \textbf{OmicMiningAgent}};

    % Tool
    \node[toolbox, below=0.8cm of agent] (tool1)
        {\texttt{omic\_analysis()}\\
         {\scriptsize NER $\rightarrow$ dataset search $\rightarrow$ DE $\rightarrow$ enrichment}};

    % Resources
    \node[resource, left=0.8cm of tool1, yshift=-0.8cm] (db1) {TOSG\\DB};
    \node[resource, below=0.2cm of db1] (db2) {GLiNER\\NER};
    \node[resource, right=0.8cm of tool1, yshift=-0.8cm] (db3) {DESeq2\\/ R};
    \node[resource, below=0.2cm of db3] (db4) {KEGG\\Enrich};

    % State I/O
    \node[stateio, below=1.8cm of tool1] (reads)
        {\textcolor{gray!60!black}{\textbf{R:}} \texttt{query}};
    \node[stateio, below=0.3cm of reads]  (writes)
        {\textcolor{blue!50!black}{\textbf{W:}} \texttt{top\_genes[]}, \texttt{disease},\\
         \quad\; \texttt{cell\_type}, \texttt{pathways[]}};

    % Connections
    \draw[conn] (agent) -- (tool1);
    \draw[conn, draw=gray!30] (tool1.south) -- ++(0,-0.3) -| (db1.north);
    \draw[conn, draw=gray!30] (tool1.south) -- ++(0,-0.3) -| (db2.north);
    \draw[conn, draw=gray!30] (tool1.south) -- ++(0,-0.3) -| (db3.north);
    \draw[conn, draw=gray!30] (tool1.south) -- ++(0,-0.3) -| (db4.north);

    \node[iolabel, left=0.1cm of reads.north west, anchor=south east] {AgentState};
    \draw[stateconn] (tool1.south east) -- ++(0.6,0) |- (reads.east);
    \draw[stateconn] (tool1.south east) -- ++(0.6,0) |- (writes.east);
\end{tikzpicture}%
}
\par\smallskip{\small (a) OmicMiningAgent}
\end{minipage}%
\hfill
% =======================  (b) BioMarkerKGAgent  =======================
\begin{minipage}[t]{0.48\textwidth}
\centering
\resizebox{\textwidth}{!}{%
\begin{tikzpicture}[node distance=0.6cm, font=\sffamily]
    % Agent header
    \node[agentmain=agentCyan] (agent)
        {\iconGraph{cyan!60!black} \enskip \textbf{BioMarkerKGAgent}};

    % Tool
    \node[toolbox, below=0.8cm of agent] (tool1)
        {\texttt{query\_knowledge\_graph()}\\
         {\scriptsize GRetriever: GAT + LLM over PrimeKG}};

    % Resources
    \node[resource, left=0.8cm of tool1, yshift=-0.8cm] (db1) {Neo4j\\PrimeKG};
    \node[resource, below=0.2cm of db1] (db2) {GRetriever\\(PyG)};
    \node[resource, right=0.8cm of tool1, yshift=-0.8cm] (db3) {OpenAI\\Embed};
    \node[resource, below=0.2cm of db3] (db4) {PCST\\Subgraph};

    % State I/O
    \node[stateio, below=1.8cm of tool1] (reads)
        {\textcolor{gray!60!black}{\textbf{R:}} \texttt{top\_genes},
         \texttt{previous\_results}};
    \node[stateio, below=0.3cm of reads]  (writes)
        {\textcolor{blue!50!black}{\textbf{W:}} \texttt{agent\_outputs[]}\\
         {\scriptsize (pathway chains, GO terms, PPIs)}};

    % Connections
    \draw[conn] (agent) -- (tool1);
    \draw[conn, draw=gray!30] (tool1.south) -- ++(0,-0.3) -| (db1.north);
    \draw[conn, draw=gray!30] (tool1.south) -- ++(0,-0.3) -| (db2.north);
    \draw[conn, draw=gray!30] (tool1.south) -- ++(0,-0.3) -| (db3.north);
    \draw[conn, draw=gray!30] (tool1.south) -- ++(0,-0.3) -| (db4.north);

    \node[iolabel, left=0.1cm of reads.north west, anchor=south east] {AgentState};
    \draw[stateconn] (tool1.south east) -- ++(0.6,0) |- (reads.east);
    \draw[stateconn] (tool1.south east) -- ++(0.6,0) |- (writes.east);
\end{tikzpicture}%
}
\par\smallskip{\small (b) BioMarkerKGAgent}
\end{minipage}

\vspace{0.8cm}

% =======================  (c) PubMedResearcher  =======================
\begin{minipage}[t]{0.48\textwidth}
\centering
\resizebox{\textwidth}{!}{%
\begin{tikzpicture}[node distance=0.6cm, font=\sffamily]
    % Agent header
    \node[agentmain=agentPurple] (agent)
        {\iconSearch{violet} \enskip \textbf{PubMedResearcher}};

    % Tools (2)
    \node[toolbox, below=0.8cm of agent, xshift=-1.5cm] (tool1)
        {\texttt{search\_pubmed()}\\
         {\scriptsize Full-text PDF/XML download}};
    \node[toolbox, below=0.8cm of agent, xshift=1.5cm] (tool2)
        {\texttt{pubmed\_search\_fast()}\\
         {\scriptsize NCBI E-utilities (abstracts)}};

    % Resources
    \node[resource, below=1.0cm of tool1] (db1) {PubMed\\API};
    \node[resource, right=0.3cm of db1] (db2) {PMC\\PDFs};
    \node[resource, below=1.0cm of tool2] (db3) {DOI\\Cache};

    % State I/O
    \node[stateio, below=2.4cm of agent] (reads)
        {\textcolor{gray!60!black}{\textbf{R:}} \texttt{top\_genes},
         \texttt{previous\_results}};
    \node[stateio, below=0.3cm of reads]  (writes)
        {\textcolor{blue!50!black}{\textbf{W:}} \texttt{paper\_dois[]},
         \texttt{paper\_pmids[]}\\
         {\scriptsize (per-gene literature with citations)}};

    % Connections
    \draw[conn] (agent.south) -- ++(0,-0.3) -| (tool1.north);
    \draw[conn] (agent.south) -- ++(0,-0.3) -| (tool2.north);
    \draw[conn, draw=gray!30] (tool1) -- (db1);
    \draw[conn, draw=gray!30] (tool1) -- (db2);
    \draw[conn, draw=gray!30] (tool2) -- (db3);

    \node[iolabel, left=0.1cm of reads.north west, anchor=south east] {AgentState};
    \draw[stateconn] (reads.east) -- ++(0.6,0) |- ($(agent.south east)+(0.6,-0.3)$);
    \draw[stateconn] (writes.east) -- ++(0.6,0) |- ($(agent.south east)+(0.6,-0.3)$);
\end{tikzpicture}%
}
\par\smallskip{\small (c) PubMedResearcher}
\end{minipage}%
\hfill
% =======================  (d) GoogleSearcher  =======================
\begin{minipage}[t]{0.48\textwidth}
\centering
\resizebox{\textwidth}{!}{%
\begin{tikzpicture}[node distance=0.6cm, font=\sffamily]
    % Agent header
    \node[agentmain=googleGray] (agent)
        {\iconSearch{gray!60!black} \enskip \textbf{GoogleSearcher}};

    % Tool
    \node[toolbox, below=0.8cm of agent] (tool1)
        {\texttt{google\_search()}\\
         {\scriptsize Custom Search API + w3m extraction}};

    % Resources
    \node[resource, left=0.8cm of tool1, yshift=-0.8cm] (db1) {Google\\CSE};
    \node[resource, right=0.8cm of tool1, yshift=-0.8cm] (db2) {w3m\\Browser};
    \node[resource, below=0.2cm of db2] (db3) {LLM\\Filter};

    % State I/O
    \node[stateio, below=1.8cm of tool1] (reads)
        {\textcolor{gray!60!black}{\textbf{R:}} \texttt{previous\_results}
         (hypotheses)};
    \node[stateio, below=0.3cm of reads]  (writes)
        {\textcolor{blue!50!black}{\textbf{W:}} \texttt{agent\_outputs[]}\\
         {\scriptsize (clinical trials, translational viability)}};

    % Connections
    \draw[conn] (agent) -- (tool1);
    \draw[conn, draw=gray!30] (tool1.south) -- ++(0,-0.3) -| (db1.north);
    \draw[conn, draw=gray!30] (tool1.south) -- ++(0,-0.3) -| (db2.north);
    \draw[conn, draw=gray!30] (tool1.south) -- ++(0,-0.3) -| (db3.north);

    \node[iolabel, left=0.1cm of reads.north west, anchor=south east] {AgentState};
    \draw[stateconn] (tool1.south east) -- ++(0.6,0) |- (reads.east);
    \draw[stateconn] (tool1.south east) -- ++(0.6,0) |- (writes.east);
\end{tikzpicture}%
}
\par\smallskip{\small (d) GoogleSearcher}
\end{minipage}

\vspace{0.8cm}

% =======================  (e) ScientistsAgent  =======================
\begin{minipage}[t]{\textwidth}
\centering
\resizebox{0.85\textwidth}{!}{%
\begin{tikzpicture}[node distance=0.6cm, font=\sffamily]
    % Agent header (wider for compound agent)
    \node[agentmain=agentTeal, minimum width=8cm] (agent)
        {\iconBulb{teal!60!black} \enskip \textbf{ScientistsAgent}
         {\small\itshape (Multi-Expert PI)}};

    % Router
    \node[rectangle, rounded corners=4pt, draw=teal!50, thick,
          fill=teal!8, minimum width=5cm, minimum height=0.9cm,
          align=center, font=\sffamily\small,
          below=0.7cm of agent] (router)
        {\iconBrain{teal!60!black} \enskip \textbf{Expert Router (PI)}\\
         {\scriptsize Structured LLM routing $\rightarrow$ selects 2--3 experts}};

    % 4 Expert boxes in a row
    \node[toolbox, fill=yellow!8, below=0.9cm of router, xshift=-5.2cm,
          minimum width=3.6cm, minimum height=1.8cm] (exp1)
        {\textbf{GenomicsExpert}\\[2pt]
         \texttt{\scriptsize query\_genomics-}\\
         \texttt{\scriptsize expert\_kb()}\\
         {\scriptsize HyperRAG KB}};

    \node[toolbox, fill=pink!8, right=0.3cm of exp1,
          minimum width=3.6cm, minimum height=1.8cm] (exp2)
        {\textbf{NeuroscienceExpert}\\[2pt]
         \texttt{\scriptsize query\_neuroscience-}\\
         \texttt{\scriptsize expert\_kb()}\\
         {\scriptsize HyperRAG KB}};

    \node[toolbox, fill=lime!8, right=0.3cm of exp2,
          minimum width=3.6cm, minimum height=1.8cm] (exp3)
        {\textbf{LongevityBiostats}\\[2pt]
         \texttt{\scriptsize query\_longevity-}\\
         \texttt{\scriptsize biostatsexpert\_kb()}\\
         {\scriptsize HyperRAG KB}};

    \node[toolbox, fill=cyan!8, right=0.3cm of exp3,
          minimum width=3.6cm, minimum height=1.8cm] (exp4)
        {\textbf{BioinformaticsExpert}\\[2pt]
         \texttt{\scriptsize query\_bioinformatics-}\\
         \texttt{\scriptsize expert\_kb()}\\
         {\scriptsize HyperRAG KB}};

    % Synthesis
    \node[rectangle, rounded corners=4pt, draw=teal!50, thick,
          fill=teal!8, minimum width=5cm, minimum height=0.9cm,
          align=center, font=\sffamily\small,
          below=0.9cm of $(exp2.south)!0.5!(exp3.south)$] (synth)
        {\iconRobot{teal!60!black} \enskip \textbf{PI Synthesis}\\
         {\scriptsize Integrates expert perspectives $\rightarrow$ ranked hypotheses}};

    % State I/O (below synthesis)
    \node[stateio, minimum width=6cm, below=0.8cm of synth, xshift=-2.5cm] (reads)
        {\textcolor{gray!60!black}{\textbf{R:}} \texttt{ALL previous\_results},
         \texttt{top\_genes}, \texttt{paper\_dois}};
    \node[stateio, minimum width=6cm, below=0.3cm of reads]  (writes)
        {\textcolor{blue!50!black}{\textbf{W:}} \texttt{agent\_outputs[]}
         + \texttt{expert\_routing\{\}}\\
         {\scriptsize (ranked hypotheses, validation plans, expert contributions)}};

    % RAG resources
    \node[resource, below=0.15cm of exp1] (rag1) {Author\\KB};
    \node[resource, below=0.15cm of exp2] (rag2) {Author\\KB};
    \node[resource, below=0.15cm of exp3] (rag3) {Author\\KB};
    \node[resource, below=0.15cm of exp4] (rag4) {Author\\KB};

    % Connections: agent -> router
    \draw[conn] (agent) -- (router);

    % Router -> experts (fan-out)
    \draw[conn, draw=teal!40] (router.south) -- ++(0,-0.25) -| (exp1.north);
    \draw[conn, draw=teal!40] (router.south) -- ++(0,-0.25) -| (exp2.north);
    \draw[conn, draw=teal!40] (router.south) -- ++(0,-0.25) -| (exp3.north);
    \draw[conn, draw=teal!40] (router.south) -- ++(0,-0.25) -| (exp4.north);

    % Experts -> synthesis (fan-in)
    \draw[conn, draw=teal!40] (exp1.south) -- ++(0,-0.95) -| (synth.north west);
    \draw[conn, draw=teal!40] (exp2.south) -- ++(0,-0.95) -| ($(synth.north)!0.33!(synth.north west)$);
    \draw[conn, draw=teal!40] (exp3.south) -- ++(0,-0.95) -| ($(synth.north)!0.33!(synth.north east)$);
    \draw[conn, draw=teal!40] (exp4.south) -- ++(0,-0.95) -| (synth.north east);

    % State connections
    \node[iolabel, left=0.1cm of reads.north west, anchor=south east] {AgentState};
    \draw[stateconn] (synth.south east) -- ++(0.6,0) |- (reads.east);
    \draw[stateconn] (synth.south east) -- ++(0.6,0) |- (writes.east);
\end{tikzpicture}%
}
\par\smallskip{\small (e) ScientistsAgent}
\end{minipage}

\caption{Close-up views of each sub-agent in the OmniCellAgent system, showing
tools, external resources, and AgentState read/write interfaces.
\textbf{(a)~OmicMiningAgent} performs differential expression analysis via a
single \texttt{omic\_analysis} tool backed by the OmniCellTOSG single-cell
database, GLiNER NER for entity extraction, and R/DESeq2 for statistical
testing; it writes the foundational \texttt{top\_genes} list that seeds all
downstream agents.
\textbf{(b)~BioMarkerKGAgent} queries the PrimeKG knowledge graph through a
GRetriever model (GAT + LLM on Neo4j) to map genes to pathways, protein
interactions, and GO terms, contextualising the omics hits within known
biology.
\textbf{(c)~PubMedResearcher} has two tools---a full-text downloader
(\texttt{search\_pubmed}) and a fast abstract search
(\texttt{pubmed\_search\_fast})---producing per-gene literature reviews with
DOI/PMID citations that populate the shared state for traceability.
\textbf{(d)~GoogleSearcher} uses the Google Custom Search API with w3m-based
page extraction to assess clinical trial pipelines and translational viability
of the proposed targets.
\textbf{(e)~ScientistsAgent} is a compound multi-expert agent: a PI-level
router selects 2--3 domain experts (Genomics, Neuroscience,
Longevity/Biostatistics, Bioinformatics), each optionally backed by a
HyperRAG knowledge base built from curated author publications; a synthesis
node integrates their perspectives into ranked mechanistic hypotheses with
proposed validation experiments.}
\label{fig:agent-closeups}
\end{figure}
